% !TeX root = ../../Book4.tex
% 纲要标题：### 18.3 导出范畴中的有界性
% 纲要标签：b4:sec:1703
% 纲要位置：计划纲要.md，第 3696 行。

\section{导出范畴中的有界性}
\label{b4:sec:1703}

复形的项有界与其同调有界并非同一条件。典范截断允许删去正合尾部，却可能失去投射或内射性，因此不能把有界对象与有界特殊消解混同。


% 纲要内容（保留纲要核验项目）：
%
% * D^b、D^+ 与 D^-
% * 项有界和同调有界的区别
% * 模范畴中的基本实例
%

\subsection{\texorpdfstring{\(D^b\)、\(D^+\) 与 \(D^-\)}{Db、D+ 与 D−}}
\label{b4:subsec:1703:01}

\begin{definition}\label{b4:def:bounded-derived}
设 \(\mathcal A\) 为局部小交换范畴，并固定局部化的集合大小约定。分别将 \(K^b(\mathcal A),K^+(\mathcal A),K^-(\mathcal A)\) 中的拟同构倒转，得到\term[youjie-daochu-fanchou]{有界导出范畴}[bounded derived category] \(D^b(\mathcal A)\)、有下界导出范畴 \(D^+(\mathcal A)\) 和有上界导出范畴 \(D^-(\mathcal A)\)。有界在此先指复形的项。
\end{definition}
\symidx{\(D^b(\mathcal A),D^+(\mathcal A),D^-(\mathcal A)\)：有界、有下界与有上界导出范畴}

\begin{proposition}\label{b4:prop:bounded-derived-full}
设 \(\mathcal A\) 为局部小交换范畴，且相关局部化均采用同一集合大小约定。自然函子 \(D^b(\mathcal A),D^+(\mathcal A),D^-(\mathcal A)\to D(\mathcal A)\) 全忠实；它们的本质像分别由上同调两端有界、有下界、有上界的对象组成。三者各为三角子范畴。
\end{proposition}
\begin{proof}
截去尾部时必须保留截口上的余圈与余边条件。例如，\(\mathbb Z\xrightarrow{2}\mathbb Z\xrightarrow{q}\mathbb Z/2\mathbb Z\) 位于次数零、一、二，原复形零调；若直接删去次数大于一的项，余下的 \(\mathbb Z\xrightarrow{2}\mathbb Z\) 却有 \(H^1=\mathbb Z/2\mathbb Z\)。这是因为末项的全部元素都变成余圈。把次数一的项改为 \(\ker q=2\mathbb Z\)，所得 \(\mathbb Z\xrightarrow{2}2\mathbb Z\) 就仍零调。

一般地，复形 \(C\) 的 \(\tau_{\le n}C\) 在 \(n\) 次取 \(\ker \mathrm{d}_C^n\)，较低次数不变，较高次数取零；它自然映入 \(C\)，并保留次数不大于 \(n\) 的上同调。\(\tau_{\ge m}C\) 在 \(m\) 次取 \(\operatorname{coker}\mathrm{d}_C^{m-1}\)，较高次数不变，较低次数取零；\(C\) 自然映到它，并保留次数不小于 \(m\) 的上同调，其余次数的上同调为零。右端取核，是为了保留原来成为余圈的条件；左端取余核，是为了仍将来自被删项的余边模掉。这些断言直接从核、像的定义核对，因而两种截断都保持拟同构。链映射在截口分别限制到核、下降到余核，给出截断后的链映射，并保持恒等与复合。

为了在同伦范畴中截断屋顶及其共同细化，还必须保证同伦的映射经截断后仍同伦。截口取核或余核依赖整个复形，不能仅由逐项加性函子保持同伦的结论代替这一步。设 \(f,g:C\to E\) 及 \(h^i:C^i\to E^{i-1}\) 满足
\[
 f^i-g^i=\mathrm{d}_E^{i-1}h^i+h^{i+1}\mathrm{d}_C^i.
\]
对 \(\tau_{\le n}\)，在 \(i<n\) 保留 \(h^i\)，在 \(i=n\) 将 \(h^n\) 限制到 \(\ker\mathrm{d}_C^n\)，在更高次数取零。截口上 \(h^{n+1}\mathrm{d}_C^n\) 消失，剩下的 \(\mathrm{d}_E^{n-1}h^n\) 取值于 \(\ker\mathrm{d}_E^n\)，所以同伦等式仍成立；在 \(n-1\) 次，限制后的 \(h^n\) 与诱导微分的复合仍为 \(h^n\mathrm{d}_C^{n-1}\)。对 \(\tau_{\ge m}\)，记截口的商映射为
\[
 q_C:C^m\to\operatorname{coker}\mathrm{d}_C^{m-1},
 \qquad q_E:E^m\to\operatorname{coker}\mathrm{d}_E^{m-1}.
\]
在 \(i>m+1\) 保留 \(h^i\)，在 \(i=m+1\) 取 \(q_Eh^{m+1}\)，在 \(i\le m\) 取零。由于 \(q_E\mathrm{d}_E^{m-1}=0\)，截口上的等式成为
\[
 (\tau_{\ge m}f-\tau_{\ge m}g)^m q_C
 =q_Eh^{m+1}\mathrm{d}_C^m,
\]
正是余核上所需的同伦等式；在 \(m+1\) 次，余核诱导的微分与 \(q_Eh^{m+1}\) 复合仍为 \(\mathrm{d}_E^mh^{m+1}\)。因此两种截断都给出 \(K(\mathcal A)\) 上的函子。

若 \(X,Y\) 的项在次数 \(n\) 以上为零，则屋顶 \(X\leftarrow U\to Y\) 可前复合 \(\tau_{\le n}U\to U\)。这里使用左屋顶，是因为截断的自然映射指向 \(U\)。因 \(H^{>n}(U)=0\)，这是拟同构，所以得到有上界的屋顶。若两个有上界的屋顶在 \(D(\mathcal A)\) 中相等，可增大 \(n\)，使这两个屋顶的全部对象在 \(n\) 以上为零，再对它们的共同细化施加 \(\tau_{\le n}\)。截断保持拟同构及 \(K(\mathcal A)\) 中的等式，且不改变两个原屋顶，因此得到有上界的共同细化。这同时证明 \(D^-\) 的满性和忠实性。对 \(D^+\)，截断的自然映射方向为 \(U\to\tau_{\ge m}U\)，因此使用分母在右的屋顶 \(X\to U\leftarrow Y\)，并把两个箭头后复合到截断对象。比较两个屋顶时，取它们全部对象的公共下界 \(m\)，截断共同细化，得到相同的结论。

若 \(X,Y\) 都只在 \([m,n]\) 内有项，先把左屋顶截到 \(\tau_{\le n}U\)，再取它的 \(\tau_{\ge m}\)。指向 \(X,Y\) 的复形映射在次数 \(m\) 都消去来自 \(m-1\) 的像，因此均通过该余核下降，得到有界屋顶；其左边仍为拟同构。比较两个有界屋顶时，选取包含它们全部对象的非零项的公共区间 \([m,n]\)，令 \(T=\tau_{\ge m}\tau_{\le n}\)。若共同细化中的等式为 \(s_1a=s_2b\)、\(f_1a=f_2b\)，施加 \(T\) 后仍有
\[
 T(s_1)T(a)=T(s_2)T(b),
 \qquad T(f_1)T(a)=T(f_2)T(b).
\]
这里 \(T\) 不改变两个原屋顶，而 \(T(a),T(b)\) 仍为拟同构，故所得细化属于 \(K^b(\mathcal A)\)。这证明 \(D^b\) 也全忠实。同样的截断论证在单向有界的屋顶范畴内直接进行，还给出 \(D^b(\mathcal A)\to D^+(\mathcal A),D^-(\mathcal A)\) 的全忠实性。

若 \(H^i(X)=0\) 对所有 \(i\notin[m,n]\) 成立，则有拟同构之字形
\[
 X\xleftarrow{\simeq}\tau_{\le n}X
 \xrightarrow{\simeq}\tau_{\ge m}\tau_{\le n}X.
\]
最右的复形只在 \([m,n]\) 内有项，故 \(X\) 在导出范畴中与项有界的复形同构。只有上界或下界时，分别用 \(\tau_{\le n}X\to X\) 或 \(X\to\tau_{\ge m}X\) 即可。反向由上同调保持拟同构显然。平移和锥保持相应的上同调有界性，故这些满子范畴对好三角封闭。
\end{proof}

\subsection{项有界与同调有界}
\label{b4:subsec:1703:02}

“项有界”是某个具体代表的性质，“上同调有界”是导出范畴中对象的性质。上命题说明后者允许更换成项有界的代表，却没有说明该代表还能同时逐项投射或逐项内射。为获得特殊项，可能必须重新加上一段无限长但正合的尾部。

\ProseParagraphBreak
例如投射消解 \(\cdots\to P_2\to P_1\to P_0\to M\to0\)，删去增广项后，把 \(P_i\) 放在次数 \(-i\)，得到复形 \(P\)。它只有 \(H^0(P)=M\)，因此代表 \(D^b\) 的对象；但它的项可能向负次数无限延伸。把正合尾部简单截掉通常会新增上同调；正确截断需要把端点项换成核或余核，而后者未必仍投射。

\ProseParagraphBreak
这也解释了为什么两个有界条件在导出函子公式中各有用处：上同调有界用于识别对象属于哪一类导出范畴，项及其投射、内射性质则用于保证逐项计算有效。证明中不可把二者无说明地互换。

\subsection{模范畴中的基本实例}
\label{b4:subsec:1703:03}

对域 \(k\)，每个向量空间复形都能非典范地分解为零微分的上同调复形与可缩复形的直和。逐次选 \(B^n\subseteq Z^n\subseteq C^n\) 的补空间，则 \(C^n\cong B^n\oplus H^n\oplus B^{n+1}\)，微分只在最后一项到下一次数第一项之间为恒等。于是 \(K(k\text{-}\mathbf{Mod})\to D(k\text{-}\mathbf{Mod})\) 已是等价，每个拟同构都是同伦等价。这里依赖所有向量空间短正合列可分裂。

\ProseParagraphBreak
整数上情况不同。两项复形 \(P=(\mathbb Z\xrightarrow{m}\mathbb Z)\) 位于次数负一、零，其中 \(m\ge2\)，自然映到 \(\mathbb Z/m\mathbb Z\)，是拟同构而非同伦等价。它给出屋顶
\[
 \mathbb Z/m\mathbb Z\xleftarrow{\simeq}P\longrightarrow\mathbb Z[1],
\]
右边的复形映射 \(f\) 在次数负一取恒等、在次数零取零，满足链映射条件。若 \(f\) 同伦于零，同伦只能在次数零给出 \(h^0:\mathbb Z\to\mathbb Z\)，而次数负一的同伦等式要求 \(h^0\circ m=1_{\mathbb Z}\)，即某个整数乘 \(m\) 等于一，这是不可能的。\cref{b4:prop:resolution-computes-derived-hom}将证明，源为这种有上界投射复形时，局部化不再使新的同伦类变成零；因此该屋顶表示非零导出态射。\secref{b4:sec:1704}还将严格算出 \(\operatorname{Hom}_D(\mathbb Z/m\mathbb Z,\mathbb Z[1])\cong\mathbb Z/m\mathbb Z\)，说明其倍数按模 \(m\) 区分。它在所有上同调对象之间诱导零映射，因为两端非零上同调位于不同次数。这说明“知道全部上同调对象和其上诱导的映射”仍不足以恢复导出范畴的态射。
